% =====================================================================
%  4. LA FORMALIZACIÓN EN LEAN
% =====================================================================
\chapter{La formalización en Lean}\label{sec:formalizacion}

AAA

\section{Visión general y decisiones de diseño}\label{sec:form-diseno}

AAA

\section{La sintaxis}\label{sec:form-sintaxis}

AAA

\section{El esquema de tautologías}\label{sec:form-tautologias}

AAA

\section{El sistema axiomático}\label{sec:form-sistema}

AAA

\section{La semántica de Kripke}\label{sec:form-semantica}

AAA

\begin{observacion}\label{obs:vacio}

No exigimos que \lean{World} sea no vacío, a diferencia de la definición
\ref{def:kripke}. Sobre el modelo (único) con tipo de mundos vacío toda
fórmula es válida vacuamente, lo cual es inofensivo: la validez cuantifica
sobre \emph{todos} los modelos, y sobra con que exista alguno no vacío para
que $\bot$ no sea válida. A cambio, las definiciones quedan más limpias (sin
arrastrar una prueba de no vacuidad) y los teoremas, marginalmente más
generales.
\end{observacion}

\section{El teorema de corrección}\label{sec:form-correccion}

AAA

\section{El teorema de completitud}\label{sec:form-completitud}

AAA

\subsection{Deducción a partir de hipótesis}\label{subsec:comp-ded}

AAA

\subsection{Consistencia y conjuntos maximales consistentes}\label{subsec:comp-cons}

AAA

\subsection{El lema de Lindenbaum, vía Zorn}\label{subsec:comp-lind}

AAA

\subsection{El modelo canónico y el lema de encajonado}\label{subsec:comp-mod}

AAA

\subsection{Los lemas de existencia y de verdad}\label{subsec:comp-ver}

AAA

\subsection{Completitud y adecuación}\label{subsec:comp-fin}

AAA
